\documentclass[12pt, letterpaper]{article}

%-------Packages---------
\usepackage{newpxtext,newpxmath} % Palatino-like text & math (modern)
\usepackage{bboldx}
\usepackage{mathtools}
\usepackage{tikz}
\usetikzlibrary{calc,intersections}
\usepackage{tikz-cd}
\usepackage[all,arc]{xy}
\usepackage{graphicx}

\usepackage{enumerate}
\usepackage[dvipsnames]{xcolor}
\usepackage{float}
\usepackage[left=1in,top=1in,right=1in,bottom=1in]{geometry}
\usepackage[nocheck]{fancyhdr}
\usepackage{multicol}
\usepackage{setspace}

\usepackage{dsfont}
\usepackage{mathrsfs}

\usepackage{tcolorbox}
\tcbuselibrary{theorems,skins,breakable}

\usepackage{xparse}
\usepackage{import}
\usepackage[utf8]{inputenc}

\usepackage{hyperref}
\usepackage{cleveref}
\usepackage{mdframed}
\usepackage{minted}
\setminted{
    % You can add other global options here, like:
    breaklines=true,
    autogobble,
    frame = single
}

\usepackage{xcolor, ifthen, tcolorbox}
\tcbuselibrary{theorems,skins,breakable}
% Theorem System
% The following environments are provided:
%   New Problem:        \begin{problem} ...
%   New Question Header:   \qh (then followed by subproblems)
%   Subproblem:     \begin{subproblem} ...
%   Problem Proof:  \begin{problemproof} ...
%   Proof:          \\begin{pf}{color} ...
%   Remark:         \rmk (sentence), \rmkb (block)

% Define custom colors
\definecolor{thmscol}{HTML}{F4565B} %For Problems
\definecolor{lemscol}{HTML}{FE844C} %For Lemmes
\definecolor{prpscol}{HTML}{00A7EF} %For proposition
\definecolor{corscol}{HTML}{23DAFF} %For corrolaries
\definecolor{rmkscol}{HTML}{6DEEC5} %For remarks
\definecolor{defscol}{HTML}{18C991} %For definitions
\definecolor{exmscol}{HTML}{7DB800} %For examples
\definecolor{clmscol}{HTML}{165c58} %For claims

%Change between dark(black)/light(white) mode
\newcommand{\colormode}{white}
\ifthenelse{\equal{\colormode}{white}}
      {\newcommand{\TextColor}{black}}
      {\newcommand{\TextColor}{white}}
\newcommand{\pbcolormain}{thmscol!80!\colormode}
\newcommand{\pbcolorsecond}{thmscol!38!\colormode}


% ============================
% Problem Counter
% ============================
\newcounter{subproblemcounter}
\newcounter{problemcounter}

\newcommand{\q}{
    \setcounter{subproblemcounter}{0}%
    \stepcounter{problemcounter} % Increment problem counter
    \color{\TextColor}Problem \arabic{problemcounter}: % Display the problem counter
}

\newcommand{\stepsub}{
    \stepcounter{subproblemcounter}%
    \color{\TextColor}(\alph{subproblemcounter})
}
% ============================
% Problem
% ============================
\newtcolorbox{pb}[1]{
  enhanced,
  frame hidden,
  titlerule=0mm,
  toptitle=1mm,
  bottomtitle=1mm,
  fonttitle=\bfseries\large,
  coltitle=black,
  colbacktitle=\pbcolormain,
  colback=\pbcolorsecond,
  title={\q #1},
  coltext=\TextColor
}

\newtcolorbox{spb}[1]{
  enhanced,
  frame hidden,
  titlerule=0mm,
  toptitle=1mm,
  bottomtitle=1mm,
  fonttitle=\bfseries\large,
  coltitle=black,
  colbacktitle=\pbcolormain,
  colback=\pbcolorsecond,
  title={\stepsub #1},
  coltext=\TextColor,
  attach boxed title to top left={xshift=3mm,yshift=-2mm},
  boxed title style={
    colback=\pbcolormain,
    colframe=\pbcolormain,
  }
}

\newtcolorbox{pbh}[1]{
    enhanced,
    frame hidden,
    titlerule=0mm,
    toptitle=1mm,
    bottomtitle=1mm,
    fonttitle=\bfseries\large,
    coltitle=black,
    colbacktitle=\pbcolormain,
    colback=\pbcolorsecond,
    title={\q #1},
    boxrule=0pt,
    interior hidden,
    attach boxed title to top left={xshift=3mm,yshift=-2mm},
    boxed title style={
        colback=\pbcolormain,
        colframe=\pbcolormain,
    }
}

\newenvironment{problem}[1]{%
  \newpage
  \begin{pb}{#1}
    }{
  \end{pb}
}

\newenvironment{subproblem}[1]{%
  \begin{spb}{#1}
    }{
  \end{spb}
}

\newenvironment{problemproof}{%
    \begin{pf}{\pbcolormain}
        }{
    \end{pf}
}

\newcommand{\qh}[1][]{
    \newpage
    \begin{pbh}{#1}\end{pbh} % Display the problem counter
}

% ============================
% Claim
% ============================

\newtcbtheorem[number within=problemcounter]{claim}{Claim}
{
    enhanced,
    frame hidden,
    titlerule=0mm,
    toptitle=1mm,
    bottomtitle=1mm,
    fonttitle=\bfseries\large,
    coltitle=black,
    colbacktitle=clmscol!50!\colormode,
    colback=clmscol!20!\colormode,
}{clm}

\NewDocumentCommand{\clm}{m+m}{
    \begin{claim*}{#1}{}
        #2
    \end{claim*}
}

\newenvironment{clmpf}{
	{\noindent{\it \textbf{Proof.}}
	\setlength{\parindent}{0pt}
	}
	\tcolorbox[blanker,breakable,left=5mm,parbox=false,
    before upper={\parindent15pt},
    after skip=10pt,
	borderline west={1mm}{0pt}{clmscol!50!\colormode}]
}{
    \textcolor{clmscol!50!\colormode}{\hbox{}\nobreak\hfill$\blacksquare$}
    \endtcolorbox
}

\NewDocumentCommand{\clmp}{m+m+m}{
    \begin{claim*}{#1}{}
        #2
    \end{claim*}

    \begin{clmpf}
        #3
    \end{clmpf}
}


% ============================
% Proof
% ============================

\newenvironment{pf}[1]{%
  \def\pfcolor{#1}% Store the color in a macro
  \noindent{\it \textbf{Proof.}}%
  \tcolorbox[blanker,breakable,left=5mm,parbox=false,%
    before upper={\parindent15pt},%
    after skip=10pt,%
    borderline west={1mm}{0pt}{\pfcolor},%
    coltext=\TextColor
  ]%
  \setlength{\parindent}{0pt}%
}{%
  \textcolor{\pfcolor}{\hbox{}\nobreak\hfill$\blacksquare$}%
  \endtcolorbox%
}

\newtcolorbox{solution}{
  blanker,
  breakable,
  left=5mm,
  parbox=false,%
  before upper={\parindent15pt},%
  after skip=10pt,%
  borderline west={1mm}{0pt}{\pbcolormain},%
  coltext=\TextColor
}


% ============================
% Remark
% ============================


\NewDocumentCommand{\rmk}{+m}{
    {\it \color{rmkscol!100!white}#1}
}

\newenvironment{remark}{
    \par
    \vspace{5pt}
    \begin{minipage}{\textwidth}
        {\par\noindent{\textbf{Remark.}}}
        \tcolorbox[blanker,breakable,left=5mm,
        before skip=10pt,after skip=10pt,
        borderline west={1mm}{0pt}{rmkscol!80!\colormode},
        coltext=\TextColor
        ]
}{
        \endtcolorbox
    \end{minipage}
    \vspace{5pt}
}

\NewDocumentCommand{\rmkb}{+m}{
    \begin{remark}
        #1
    \end{remark}
}



\newcounter{subsubproblemcounter}
\newcommand{\stepsubsub}{
    \stepcounter{subsubproblemcounter}%
    \color{\TextColor}\roman{subsubproblemcounter}.
}
\renewcommand{\q}{
    \setcounter{subproblemcounter}{0}%
    \setcounter{subsubproblemcounter}{0}%
    \stepcounter{problemcounter} % Increment problem counter
    \color{\TextColor}Problem \arabic{problemcounter}: % Display the problem counter
}
\newtcolorbox{sspb}[1]{
  enhanced,
  frame hidden,
  titlerule=0mm,
  toptitle=1mm,
  bottomtitle=1mm,
  fonttitle=\bfseries\large,
  coltitle=black,
  colbacktitle=\pbcolormain,
  colback=\pbcolorsecond,
  title={\stepsubsub #1},
  coltext=\TextColor,
  attach boxed title to top left={xshift=3mm,yshift=-2mm},
  boxed title style={
    colback=\pbcolormain,
    colframe=\pbcolormain,
  }
}
\newenvironment{subsubproblem}[1]{%
  \begin{sspb}{#1}
    }{
  \end{sspb}
}

% Define a new command for problems
\renewcommand{\headrule}{\color{\TextColor}\hrulefill}
\newcommand{\C}{\mathbb{C}}
\newcommand{\R}{\mathbb{R}}
\newcommand{\Q}{\mathbb{Q}}
\newcommand{\Z}{\mathbb{Z}}
\newcommand{\N}{\mathbb{N}}
\renewcommand{\P}{\mathbb{P}}
\newcommand{\F}{\mathbb{F}}
\newcommand{\contr}{\Rightarrow \Leftarrow}
\newcommand{\s}{\(\,\)}
\renewcommand{\t}[1]{\text{#1}}


\newcommand{\bit}{{\{0,1\}}}
\newcommand{\from}{\leftarrow}
\newcommand{\com}{\mathrm{com}}
\newcommand{\sample}{\overset{\$}{\leftarrow}}

% Meta data
\setstretch{1.2}

\begin{document}
\color{\TextColor}
\pagecolor{\colormode}
\setlength{\parindent}{0pt}
\pagestyle{fancy}
\fancyhf{}
\fancyhead[LOE]{\color{\TextColor}\slshape{\textbf{Vincent Ferrara}}}
\fancyhead[ROE]{\color{\TextColor}\jobname --- \thepage}
\fancyhead[C]{\color{\TextColor}EECS 484}

\begin{problem}{}
    Determine if the following statements describe a conceptual/logical, physical, or external
representation/schema of data.
\end{problem}
\begin{subproblem}{}
    Employers can access an employee's name, job title, salary, and work schedule, but not their
personal home address and bank account information.
\end{subproblem}
\begin{solution}
    c. external
\end{solution}

\begin{subproblem}{}
    All data is stored in tables sorted by category represented as bytes within a server hosted by
AWS.
\end{subproblem}
\begin{solution}
    b. physical
\end{solution}

\begin{subproblem}{}
    The table “Students” has ID, Name, Date of Birth, Height, and Address fields. ID is unique for
each item.
\end{subproblem}
\begin{solution}
    a. conceptual/logical
\end{solution}


\qh
\begin{subproblem}{}
    For the given instances (i.e. snapshot/subset) of a student relation, identify the following attributes
(set) as \underline{might be a key} or \underline{definitely not a key} for the \textbf{entire} student relation, and provide
explanations.
\end{subproblem}
\newtcbox{\inlinecode}{on line, colback=gray!20, colframe=gray!20, boxrule=0pt, arc=2pt, boxsep=1pt, left=2pt, right=2pt, top=1pt, bottom=1pt}
\begin{solution}
    \begin{enumerate}
        \item \inlinecode{gpa} \underline{might be a key} since all gpa values are distinct in the table given.
        \item \inlinecode{sid} \underline{might be a key} since all id's are distinct in the table given.
        \item \inlinecode{first name} is \underline{definitely not a key} since there are two smith's.
        \item \inlinecode{\{first name, contact\}} \underline{might be a key} since the combination of these two attributes is unique across all rows in the table.
    \end{enumerate}
\end{solution}

{
\renewcommand{\stepsub}{
    \stepcounter{subproblemcounter}%
    \color{\TextColor}\arabic{subproblemcounter}.
}
\setcounter{subproblemcounter}{4}
\begin{subproblem}{}
    List all single attributes that are \underline{\textbf{definitely not} candidate keys}. Why are these attributes
definitely not candidate keys? If all attributes could be candidate keys, explain why.
\end{subproblem}
\begin{solution}
    \inlinecode{first name} and \inlinecode{age} are definitely not candidate keys because there are duplicates in the snapshot. A candidate
    key must uniquely identify every tuple.
\end{solution}

\begin{subproblem}{}
    Which single attribute is \underline{\textbf{definitely } a candidate key}? Why? If there are no candidate keys,
explain why.
\end{subproblem}
\begin{solution}
    While \inlinecode{gpa} and \inlinecode{contact} are also unique in this snapshot, \inlinecode{sid} is the only single attribute that is definitely a candidate key for the entire student relation. An ID is designed to be a unique identifider.
\end{solution}}

\qh
\begin{subproblem}{}
    A nurse, Alice (eid=30), in a given shift, 12pm-8pm (sid=3), can care for more than one
patient.
\end{subproblem}
\begin{solution}{}
    True
\end{solution}

\begin{subproblem}{}
    The Nurse entity does not have a primary key.
\end{subproblem}
\begin{solution}
    False
\end{solution}

\begin{subproblem}{}
    A patient, Dave (pid=23) cannot be “cared for” during multiple shifts by nurse Alice
(eid=30).
\end{subproblem}
\begin{solution}
    True
\end{solution}

\begin{subproblem}{}
    This design enforces that every patient is taken care of in each shift.
\end{subproblem}
\begin{solution}
    IDK come back True
\end{solution}

\begin{subproblem}{}
    This design enforces that in each shift, at least one nurse is taking care of some patients.
\end{subproblem}
\begin{solution}
    True
\end{solution}

\begin{subproblem}{}
    A patient, Dave (pid=23), can be taken care of by two different nurses in two different shifts.
\end{subproblem}
\begin{solution}
    False
\end{solution}

\begin{subproblem}{}
    This design allows a patient that is not taken care of by any nurse.
\end{subproblem}
\begin{solution}
    False
\end{solution}

\qh
\begin{subproblem}{}
    Does this design allow multiple musicians to collaborate on an album?
\end{subproblem}
\begin{solution}
    No
\end{solution}

\begin{subproblem}{}
    Does this design allow multiple musicians to give concerts in the same place and date?
\end{subproblem}
\begin{solution}
    Yes
\end{solution}

\begin{subproblem}{}
    Does this design allow musicians to collaborate with other musicians to perform a song?
Musicians collaborate when they perform the same song, in the same place, on the same
date.
\end{subproblem}
\begin{solution}
    Yes
\end{solution}

\begin{subproblem}{}
    Does this design allow a musician to give multiple concerts in the same place on different
dates?
\end{subproblem}
\begin{solution}
    Yes
\end{solution}

\begin{subproblem}{}
    Does this design allow a musician to perform songs produced by another musician, in a
concert?
\end{subproblem}
\begin{solution}
    Yes
\end{solution}

\begin{subproblem}{}
    Does this design allow a musician to give a concert without performing any songs?
\end{subproblem}
\begin{solution}
    No
\end{solution}

\begin{subproblem}{}
    Does this design allow for a song that is not in any album?
\end{subproblem}
\begin{solution}
    Yes
\end{solution}

\begin{subproblem}{}
    Does this design allow a musician to produce an album that contains no songs?
\end{subproblem}
\begin{solution}
    Yes
\end{solution}

\begin{subproblem}{}
    Does this design allow a musician to produce many albums?
\end{subproblem}
\begin{solution}
    Yes
\end{solution}

\begin{subproblem}{}
    Does this design allow for a musician who only sings (i.e., does not play an instrument)?
\end{subproblem}
\begin{solution}
    No
\end{solution}

\qh
\begin{subproblem}{}
    The fewest number of tables that can be used to represent this ER diagram without
redundancy is 2
\end{subproblem}
\begin{solution}
    True
\end{solution}

\begin{subproblem}{}
    In the least redundant design, 'mid' is a foreign key in the Employee table
\end{subproblem}
\begin{solution}
    True
\end{solution}

\begin{subproblem}{}
    In the least redundant design, 'mid' is a foreign key in the Employee table
\end{subproblem}
\begin{solution}
    True
\end{solution}

\begin{subproblem}{}
    In the least redundant design, ‘mid’ is a foreign key in the Employee table
\end{subproblem}
\begin{solution}
    False
\end{solution}

\begin{subproblem}{}
    In the least redundant design, ‘mid’ is a foreign key in the Employee table
\end{subproblem}
\begin{solution}
    False
\end{solution}

\begin{subproblem}{}
    Create a table for \textbf{Employee} by SQL statement, in the least redundant design, including all the
necessary integrity constraints shown by the above ER diagram
\end{subproblem}
\begin{solution}
    \begin{minted}{sql}
        CREATE TABLE Employee (
            eid INTEGER PRIMARY KEY,
            name CHAR(20),
            age INTEGER,
            mid INTEGER,
            FOREIGN KEY (mid) REFERENCES Manager
        );
    \end{minted}
\end{solution}
\end{document}
